\documentclass[a4paper]{report}
% basics
\usepackage[utf8]{inputenc}
\usepackage[T1]{fontenc}
\usepackage{comment}
\usepackage{textcomp}
\usepackage[hyphens]{url}
\usepackage[english]{babel}
\usepackage[style=alphabetic,maxcitenames=1]{biblatex}
\usepackage[colorlinks=true,linkcolor=cyan,urlcolor=magenta,citecolor=violet]{hyperref}
\usepackage{graphicx}
\usepackage{float}
\usepackage{booktabs}
\usepackage{emptypage}
\usepackage{subcaption}
\usepackage{lipsum}
\usepackage[usenames,dvipsnames]{xcolor}
% quiver style
\usepackage{tikz-cd}
% `calc` is necessary to draw curved arrows.
\usetikzlibrary{calc}
% `pathmorphing` is necessary to draw squiggly arrows.
\usetikzlibrary{decorations.pathmorphing}

% A TikZ style for curved arrows of a fixed height, due to AndréC.
\tikzset{curve/.style={settings={#1},to path={(\tikztostart)
					.. controls ($(\tikztostart)!\pv{pos}!(\tikztotarget)!\pv{height}!270:(\tikztotarget)$)
					and ($(\tikztostart)!1-\pv{pos}!(\tikztotarget)!\pv{height}!270:(\tikztotarget)$)
					.. (\tikztotarget)\tikztonodes}},
	settings/.code={\tikzset{quiver/.cd,#1}
			\def\pv##1{\pgfkeysvalueof{/tikz/quiver/##1}}},
	quiver/.cd,pos/.initial=0.35,height/.initial=0}

% TikZ arrowhead/tail styles.
\tikzset{tail reversed/.code={\pgfsetarrowsstart{tikzcd to}}}
\tikzset{2tail/.code={\pgfsetarrowsstart{Implies[reversed]}}}
\tikzset{2tail reversed/.code={\pgfsetarrowsstart{Implies}}}
% TikZ arrow styles.
\tikzset{no body/.style={/tikz/dash pattern=on 0 off 1mm}}

% useful macro for class
\newcommand{\at}[3]{\left.#1\right\vert_{#2}^{#3}}
\newcommand\quotient[2]{
	\mathchoice
	{% \displaystyle
		\text{\raise1ex\hbox{$#1$}\Big/\lower1ex\hbox{$#2$}}%
	}
	{% \textstyle
		#1\,/\,#2
	}
	{% \scriptstyle
		#1\,/\,#2
	}
	{% \scriptscriptstyle  
		#1\,/\,#2
	}
}

\newcommand{\True}{\textsf{True}}
\newcommand{\T}{\textsf{T}}
\newcommand{\False}{\textsf{False}}
\newcommand{\F}{\textsf{F}}
\newcommand{\act}{\rotatebox[origin=c]{-180}{\(\,\circlearrowright\,\)}}

\usepackage{physics}
\usepackage{complexity}

\DeclareMathOperator{\id}{id}
\DeclareMathOperator{\Span}{span}
\DeclareMathOperator{\supp}{supp}
\DeclareMathOperator{\vol}{vol}
\DeclareMathOperator{\dist}{dist}
\DeclareMathOperator{\diam}{diam}
\DeclareMathOperator{\im}{Im}
\DeclareMathOperator{\sgn}{sgn}
\DeclareMathOperator{\Int}{Int}
\DeclareMathOperator{\diag}{diag}
\DeclareMathOperator{\dom}{dom}
\DeclareMathOperator*{\argmax}{arg\,max}
\DeclareMathOperator*{\argmin}{arg\,min}
\DeclareMathOperator{\Var}{Var}
\DeclareMathOperator{\Cov}{Cov}

\let\implies\Rightarrow
\let\impliedby\Leftarrow
\let\iff\Leftrightarrow

\usepackage{stmaryrd} % for \lightning
\newcommand\conta{\scalebox{1.1}{\(\lightning\)}}

\usepackage{bm}
\usepackage{bbm}

% figure support
\usepackage{import}
\usepackage{xifthen}
\pdfminorversion=7
\usepackage{pdfpages}
\usepackage{transparent}
\newcommand{\incfig}[1]{%
	\def\svgwidth{\columnwidth}
	\import{./Figures/}{#1.pdf_tex}
}


\usepackage{amsmath, amsfonts, mathtools, amsthm, amssymb}
\usepackage{geometry}
\usepackage{mathrsfs}
\usepackage{cancel}
\usepackage{systeme}

\usepackage[inline, shortlabels]{enumitem}
\usepackage{multicol}
\setlength\multicolsep{0pt}

\usepackage{caption}
\captionsetup{belowskip=0pt}
\geometry{a4paper,left=2.54cm,right=2.54cm,top=2.54cm,bottom=2.54cm}

% for the big braces
\usepackage{bigdelim}

% correct
\definecolor{correct}{HTML}{009900}
\newcommand\correct[2]{\ensuremath{\:}{\color{red}{#1}}\ensuremath{\to }{\color{correct}{#2}}\ensuremath{\:}}
\newcommand\green[1]{{\color{correct}{#1}}}

% hide parts
\newcommand\hide[1]{}

% si unitx
\usepackage{siunitx}
\sisetup{locale = UK}
% \renewcommand\vec[1]{\mathbf{#1}}
\newcommand\mat[1]{\mathbf{#1}} % Boldface matrix notation


% tikz
\usepackage{tikz}
\usetikzlibrary{intersections, angles, quotes, positioning}
\usetikzlibrary{arrows.meta}
\usepackage{pgfplots}
\pgfplotsset{compat=1.13}


\tikzset{
    force/.style={thick, {Circle[length=2pt]}-stealth, shorten <=-1pt}
}

% Algorithm Env
\usepackage[linesnumbered,lined,vlined,ruled,commentsnumbered,resetcount,algochapter]{algorithm2e}
\SetKwComment{Comment}{// }{}
\SetArgSty{textsl}
\def\algocflineautorefname{Algorithm}
\counterwithin{algocfline}{chapter}

% Clear out the item autoref name
\makeatletter
\def\itemautorefname{\@gobble}
\makeatother

% theorems
\makeatother
\usepackage{thmtools}
\usepackage[framemethod=TikZ]{mdframed}

\mdfsetup{skipabove=1em,skipbelow=0em}

\theoremstyle{definition}

\declaretheoremstyle[
    headfont=\bfseries\sffamily\color{ForestGreen!70!black}, bodyfont=\normalfont,
    mdframed={
            linewidth=2pt,
            rightline=false, topline=false, bottomline=false,
            linecolor=ForestGreen, backgroundcolor=ForestGreen!5,
            nobreak=false
        }
]{thmgreenbox}

\declaretheoremstyle[
    headfont=\bfseries\sffamily\color{ForestGreen!70!black}, bodyfont=\normalfont,
    mdframed={
            linewidth=2pt,
            rightline=false, topline=false, bottomline=false,
            linecolor=ForestGreen, backgroundcolor=ForestGreen!8,
            nobreak=false
        }
]{thmgreen2box}

\declaretheoremstyle[
    headfont=\bfseries\sffamily\color{NavyBlue!70!black}, bodyfont=\normalfont,
    mdframed={
            linewidth=2pt,
            rightline=false, topline=false, bottomline=false,
            linecolor=NavyBlue, backgroundcolor=NavyBlue!5,
            nobreak=false
        }
]{thmbluebox}

\declaretheoremstyle[
    headfont=\bfseries\sffamily\color{TealBlue!70!black}, bodyfont=\normalfont,
    mdframed={
            linewidth=2pt,
            rightline=false, topline=false, bottomline=false,
            linecolor=TealBlue,
            nobreak=false
        }
]{thmblueline}

\declaretheoremstyle[
    headfont=\bfseries\sffamily\color{RawSienna!70!black}, bodyfont=\normalfont,
    mdframed={
            linewidth=2pt,
            rightline=false, topline=false, bottomline=false,
            linecolor=RawSienna, backgroundcolor=RawSienna!5,
            nobreak=false
        }
]{thmredbox}

\declaretheoremstyle[
    headfont=\bfseries\sffamily\color{RawSienna!70!black}, bodyfont=\normalfont,
    mdframed={
            linewidth=2pt,
            rightline=false, topline=false, bottomline=false,
            linecolor=RawSienna, backgroundcolor=RawSienna!8,
            nobreak=false
        }
]{thmred2box}

\declaretheoremstyle[
    headfont=\bfseries\sffamily\color{SeaGreen!70!black}, bodyfont=\normalfont,
    mdframed={
            linewidth=2pt,
            rightline=false, topline=false, bottomline=false,
            linecolor=SeaGreen, backgroundcolor=SeaGreen!2,
            nobreak=false
        }
]{thmgreen3box}

\declaretheoremstyle[
    headfont=\bfseries\sffamily\color{WildStrawberry!70!black}, bodyfont=\normalfont,
    mdframed={
            linewidth=2pt,
            rightline=false, topline=false, bottomline=false,
            linecolor=WildStrawberry, backgroundcolor=WildStrawberry!5,
            nobreak=false
        }
]{thmpinkbox}

\declaretheoremstyle[
    headfont=\bfseries\sffamily\color{MidnightBlue!70!black}, bodyfont=\normalfont,
    mdframed={
            linewidth=2pt,
            rightline=false, topline=false, bottomline=false,
            linecolor=MidnightBlue, backgroundcolor=MidnightBlue!5,
            nobreak=false
        }
]{thmblue2box}

\declaretheoremstyle[
    headfont=\bfseries\sffamily\color{Gray!70!black}, bodyfont=\normalfont,
    mdframed={
            linewidth=2pt,
            rightline=false, topline=false, bottomline=false,
            linecolor=Gray, backgroundcolor=Gray!5,
            nobreak=false
        }
]{notgraybox}

\declaretheoremstyle[
    headfont=\bfseries\sffamily\color{Gray!70!black}, bodyfont=\normalfont,
    mdframed={
            linewidth=2pt,
            rightline=false, topline=false, bottomline=false,
            linecolor=Gray,
            nobreak=false
        }
]{notgrayline}

% \declaretheoremstyle[
% 	headfont=\bfseries\sffamily\color{RawSienna!70!black}, bodyfont=\normalfont,
% 	numbered=no,
% 	mdframed={
% 			linewidth=2pt,
% 			rightline=false, topline=false, bottomline=false,
% 			linecolor=RawSienna, backgroundcolor=RawSienna!1,
% 		},
% 	qed=\qedsymbol
% ]{thmproofbox}

\declaretheoremstyle[
    headfont=\bfseries\sffamily\color{NavyBlue!70!black}, bodyfont=\normalfont,
    numbered=no,
    mdframed={
            linewidth=2pt,
            rightline=false, topline=false, bottomline=false,
            linecolor=NavyBlue, backgroundcolor=NavyBlue!1,
            nobreak=false
        }
]{thmexplanationbox}

\declaretheoremstyle[
    headfont=\bfseries\sffamily\color{WildStrawberry!70!black}, bodyfont=\normalfont,
    numbered=no,
    mdframed={
            linewidth=2pt,
            rightline=false, topline=false, bottomline=false,
            linecolor=WildStrawberry, backgroundcolor=WildStrawberry!1,
            nobreak=false
        }
]{thmanswerbox}

\declaretheoremstyle[
    headfont=\bfseries\sffamily\color{Violet!70!black}, bodyfont=\normalfont,
    mdframed={
            linewidth=2pt,
            rightline=false, topline=false, bottomline=false,
            linecolor=Violet, backgroundcolor=Violet!1,
            nobreak=false
        }
]{conjpurplebox}

\declaretheorem[style=thmgreenbox, name=Definition, numberwithin=section]{definition}
\declaretheorem[style=thmgreen2box, name=Definition, numbered=no]{definition*}
\declaretheorem[style=thmredbox, name=Theorem, numberwithin=section]{theorem}
\declaretheorem[style=thmred2box, name=Theorem, numbered=no]{theorem*}
\declaretheorem[style=thmredbox, name=Lemma, numberwithin=section]{lemma}
\declaretheorem[style=thmredbox, name=Proposition, numberwithin=section]{proposition}
\declaretheorem[style=thmredbox, name=Corollary, numberwithin=section]{corollary}
\declaretheorem[style=thmpinkbox, name=Problem, numberwithin=section]{problem}
\declaretheorem[style=thmpinkbox, name=Problem, numbered=no]{problem*}
\declaretheorem[style=thmblue2box, name=Claim, numbered=no]{claim}
\declaretheorem[style=conjpurplebox, name=Conjecture, numberwithin=section]{conjecture}

% Redefine proof environment to get a full control.
\makeatletter
\renewenvironment{proof}[1][\proofname]{\par
    \pushQED{\qed}%
    \normalfont \topsep-2\p@\@plus6\p@\relax
    \trivlist
    \item[\hskip\labelsep
                \color{RawSienna!70!black}\sffamily\bfseries
                #1\@addpunct{.}]\ignorespaces
    \begin{mdframed}[linewidth=2pt,rightline=false, topline=false, bottomline=false,linecolor=RawSienna, backgroundcolor=RawSienna!1]
        }{%
        \popQED\endtrivlist\@endpefalse
    \end{mdframed}
}
\makeatother

\declaretheorem[style=thmbluebox, numbered=no, name=Example]{eg}
\declaretheorem[style=thmexplanationbox, numbered=no, name=Proof]{tmpexplanation}
\newenvironment{explanation}[1][]{\vspace{-10pt}\pushQED{\(\circledast\)}\begin{tmpexplanation}}{\null\hfill\popQED\end{tmpexplanation}}

\declaretheorem[style=thmblueline, numbered=no, name=Remark]{remark}
\declaretheorem[style=thmblueline, numbered=no, name=Note]{note}
\declaretheorem[style=thmpinkbox, numbered=no, name=Exercise]{exercise}
\declaretheorem[style=notgrayline, numbered=no, name=As previously seen]{prev}
\declaretheorem[style=thmgreen3box, numbered=no, name=Intuition]{intuition}
\declaretheorem[style=notgraybox, numbered=no, name=Notation]{notation}
\declaretheorem[style=thmanswerbox, numbered=no, name=Answer]{tmpanswer}
\newenvironment{answer}[1][]{\vspace{-10pt}\pushQED{\(\circledast\)}\begin{tmpanswer}}{\null\hfill\popQED\end{tmpanswer}}


\usepackage{etoolbox}
\renewcommand{\qed}{\null\hfill\(\blacksquare\)}

\makeatletter

\def\testdateparts#1{\dateparts#1\relax}
\def\dateparts#1 #2 #3 #4 #5\relax{
    \marginpar{\small\textsf{\mbox{#1 #2 #3 #5}}}
}

\def\@lecture{}%
\newcommand{\lecture}[3]{
    \ifthenelse{\isempty{#3}}{%
        \def\@lecture{Lecture #1}%
    }{%
        \def\@lecture{Lecture #1: #3}%
    }%
    \section*{\@lecture}
    \marginpar{\small\textsf{\mbox{#2}}}
}
\usepackage{pgffor}%
\newcommand{\lec}[2]{%
    \foreach \c in {#1,...,#2}{%
            \IfFileExists{Lectures/lec_\c.tex} {%
                \input{Lectures/lec_\c.tex}%
            }{}%
        }%
}

% fancy headers
\usepackage{fancyhdr}
\pagestyle{fancy}
\fancyhead[L]{}
\fancyhead[R]{\@lecture}
\fancyfoot[L]{}
\fancyfoot[R]{\thepage}
\fancyfoot[C]{\leftmark}

\makeatother

% notes
\usepackage[color=pink]{todonotes}
\usepackage{marginnote}
\let\marginpar\marginnote

% Fix some stuff
% %http://tex.stackexchange.com/questions/76273/multiple-pdfs-with-page-group-included-in-a-single-page-warning
\pdfsuppresswarningpagegroup=1

% Appendix environment
\usepackage{appendix}
\def\chapterautorefname{Section}
\def\sectionautorefname{Section}
\def\appendixautorefname{Appendix}
\renewcommand\appendixname{Appendix}
\renewcommand\appendixtocname{Appendix}
\renewcommand\appendixpagename{Appendix}
% begin appendix autoref patch [\autoref subsections in appendix](https://tex.stackexchange.com/questions/149807/autoref-subsections-in-appendix)
\makeatletter
\patchcmd{\hyper@makecurrent}{%
    \ifx\Hy@param\Hy@chapterstring
        \let\Hy@param\Hy@chapapp
    \fi
}{%
    \iftoggle{inappendix}{%true-branch
        % list the names of all sectioning counters here
        \@checkappendixparam{chapter}%
        \@checkappendixparam{section}%
        \@checkappendixparam{subsection}%
        \@checkappendixparam{subsubsection}%
        \@checkappendixparam{paragraph}%
        \@checkappendixparam{subparagraph}%
    }{}%
}{}{\errmessage{failed to patch}}

\newcommand*{\@checkappendixparam}[1]{%
    \def\@checkappendixparamtmp{#1}%
    \ifx\Hy@param\@checkappendixparamtmp
        \let\Hy@param\Hy@appendixstring
    \fi
}
\makeatletter

\newtoggle{inappendix}
\togglefalse{inappendix}

\apptocmd{\appendix}{\toggletrue{inappendix}}{}{\errmessage{failed to patch}}
\apptocmd{\subappendices}{\toggletrue{inappendix}}{}{\errmessage{failed to patch}}
% end appendix autoref patch


\makeatletter
    \def\cleardoublepage{\clearpage%
        \if@twoside
            \ifodd\c@page\else
                \vspace*{\fill}
                \hfill
                \begin{center}
                This page intentionally left blank.
                \end{center}
                \vspace{\fill}
                \thispagestyle{empty}
                \newpage
                \if@twocolumn\hbox{}\newpage\fi
            \fi
        \fi
    }
\makeatother

\setcounter{tocdepth}{3}
\author{Krit Pio Nicol}
\title{Note Template}

\thispagestyle{empty}
\addbibresource{ref.bib}

\usepackage{adjustbox}
\usepackage{centernot}



\begin{document}



%%%%%%%%%%% TITLE PAGE

\begin{titlepage}
\begin{figure}[t]
    \centering\includegraphics[width=0.5\textwidth]{Figures/usilogo.png}
\end{figure}

\begin{center}
    \textsc{\LARGE{USI Lugano\\}}
	\textsc{ \LARGE{Computer Science Department\\ }}
	%\textnormal{ \LARGE{Corso di Laurea Triennale/Magistrale in ???\\}}
	\vspace{40mm}
	\fontsize{10mm}{7mm}\selectfont 
    \textup{Computer Networking book}\\
\end{center}

\vspace{25mm}

\begin{minipage}[t]{0.47\textwidth}
	\textnormal{\large{\bf Professor\\}}
	{\large Prof. Santini\\ \\}
 \textnormal{\large{\bf TA's\\}}
	{\large M. Laporte, L. Alecci}
\end{minipage}\hfill\begin{minipage}[t]{0.47\textwidth}\raggedleft
	\textnormal{\large{\bf Student\\}}
	{\large Krit Pio Nicol}
\end{minipage}

\vspace{20mm}

\centering{\large{Year 2024 }}

\end{titlepage}

\begin{abstract}

    \begin{itemize}
        \item The Internet is the largest engineered system in history, comprising millions of computers and communication devices, billions of users, and a vast array of Internet-connected "things".
        \item Despite its size and complexity, it is possible to understand how the Internet works, thanks to guiding principles and structures.
        \item The aim is to provide a modern introduction to computer networking, emphasizing principles and practical insights for understanding current and future networks.
        \item This overview includes:
        \begin{itemize}
            \item Basic terminology and concepts in computer networking.
            \item Examination of network hardware and software components, including end systems, network applications, links, switches, access networks, and physical media.
            \item Understanding the Internet as a network of networks and how these networks interconnect.
            \item Discussion on the core aspects of computer networks, such as delay, loss, and throughput, including simple quantitative models for end-to-end throughput and delay.
            \item Key architectural principles in computer networking, including protocol layering and service models.
            \item Overview of security challenges and types of attacks in computer networks, with strategies for enhancing security.
            \item A brief history of computer networking to close the chapter.
        \end{itemize}
        \item The chapter sets the context for the rest of the book, aiming to see "the forest through the trees" by providing a broad picture without losing sight of the detailed aspects of computer networking.
    \end{itemize}
    
\end{abstract}

\newpage

\tableofcontents


%─────Formatting_References───────────────────────────────────────────────────────────────────────────────────────────────────────────────────────────
\chapter{Introduction}
\lecture{1}{13 Oct. 08:00}{First Lecture}



\section{Useful Environment}
We now see some common environment you'll need to complete your note.

\begin{definition}[Natural number]\label{def}
  We denote the set of \emph{natural numbers} as \(\mathbb{\MakeUppercase{n}} \).
\end{definition}


\begin{lemma}[Useful lemma]\label{lma}
  Given the axioms of \hyperref[def]{natural numbers \(\mathbb{\MakeUppercase{n}} \)}, we have
  \[
    0\neq 1.
  \]
\end{lemma}
\begin{proof}[An obvious proof]
  Obvious.
\end{proof}
\begin{proposition}[Useful proposition]\label{prop}
  From \autoref{lma}, we have
  \[
    0<1.
  \]
\end{proposition}
\begin{exercise}
  Prove that \(1 < 2\).
\end{exercise}
\begin{answer}
  We note the following.
  \begin{note}
    We have \autoref{prop}! We can use it iteratively!
  \end{note}
  With the help of \autoref{lma}, this holds trivially.
\end{answer}
\begin{eg}
  We now can have \(a < b\) for \(a < b\)!
\end{eg}
\begin{explanation}
  Iteratively apply the exercise we did above.
\end{explanation}
\begin{remark}
  We see that \autoref{prop} is really powerful. We now give an immediate application of it.
\end{remark}

\begin{theorem}[Mass-energy equivalence]\label{thm}
  Given \autoref{prop}, we then have
  \[
    E = mc^2.
  \]
\end{theorem}
\begin{proof}
  The blank left for me is too small,\footnote{\url{https://en.wikipedia.org/wiki/Richard_Feynman}} hence we put the proof in \autoref{appendix}.
\end{proof}

From \autoref{thm}, we then have the following.
\begin{corollary}[Riemann hypothesis]\label{col}
  The real part of every nontrivial zero of the Riemann zeta function is \(\frac{1}{2}\), where the Riemann zeta function is just
  \[
    \zeta (s)=\sum _{n=1}^{\infty }{\frac {1}{n^{s}}}={\frac {1}{1^{s}}}+{\frac {1}{2^{s}}}+{\frac {1}{3^{s}}}+\cdots.
  \]
\end{corollary}
\begin{proof}
  The proof should be trivial, we left it to you.\todo{DIY}
\end{proof}
\begin{prev}
  We see that \autoref{lma} is really helpful in the proof!
\end{prev}

\subsubsection{Internal Link}
You should see all the common usages of internal links. Additionally, we can use citations as~\cite{newton1726philosophiae}, which just link to the reference page!

\section{Figures}
A simple demo for drawing:
\begin{figure}[H]
  \centering
  \incfig{test}
  \caption[Caption]{A \(3\)-torus.\footnotemark}
  \label{fig:test}
\end{figure}
\footnotetext{For detailed information, please see \url{https://github.com/sleepymalc/VSCode-LaTeX-Inkscape}.}

\section{Commutative Diagram}
We can use the package \texttt{tikz-cd} to draw some commutative diagram.
\begin{eg}
  The cellular homology agrees with singular homology.
\end{eg}
\begin{explanation}
  The following commutative diagram shows everything.

  \adjustbox{scale=0.85,center}{%
    \begin{tikzcd}[column sep=tiny]
      &&&& {\color{red}0} \\
      & {\color{red}0} && {\color{red}H_n(X^{n+1})\cong H_n(X)} \\
      && {\color{red}H_n(X^n)} \\
      \ldots & {\color{red}H_{n+1}(X^{n+1}, X^n)} && {\color{red}H_n(X^n, X^{n-1})} && {H_{n-1}(X^{n-1}, X^{n-2})} & \ldots \\
      &&&& {\color{red}H_{n-1}(X^{n-1})} \\
      &&& 0
      \arrow[color={rgb,255:red,214;green,92;blue,92}, from=2-2, to=3-3]
      \arrow["{\partial_{n+1}}", color={rgb,255:red,214;green,92;blue,92}, from=4-2, to=3-3]
      \arrow[color={rgb,255:red,214;green,92;blue,92}, from=2-4, to=1-5]
      \arrow[color={rgb,255:red,214;green,92;blue,92}, from=3-3, to=2-4]
      \arrow[from=4-6, to=4-7]
      \arrow["{d_n}"', from=4-4, to=4-6]
      \arrow["{\partial_n}"', color={rgb,255:red,214;green,92;blue,92}, from=4-4, to=5-5]
      \arrow["{j_{n-1}}"', from=5-5, to=4-6]
      \arrow[from=6-4, to=5-5]
      \arrow["{d_{n+1}}", from=4-2, to=4-4]
      \arrow["{j_n}", color={rgb,255:red,214;green,92;blue,92}, from=3-3, to=4-4]
      \arrow[from=4-1, to=4-2]
    \end{tikzcd}
  }
\end{explanation}

\section{Fancy Stuffs}
With this header, you can achieve some cool things. For example, we can have multiple definitions under a parent environment, while maintains the numbering of definition. This is achieved by \texttt{definition*} environment with \texttt{definition} inside. For example, we can have the following.
\begin{definition*}
  We have the following number system.
  \begin{definition}[Rational number]\label{def:rational}
    The set of \emph{rational number}, denote as \(\mathbb{\MakeUppercase{q}} \).
  \end{definition}
  \begin{definition}[Real number]\label{def:real}
    The set of \emph{real number}, denote as \(\mathbb{\MakeUppercase{r}} \).
  \end{definition}
  \begin{definition}[Complex number]\label{def:complex}
    The set of \emph{complex number}, denote as \(\mathbb{\MakeUppercase{c}} \).
  \end{definition}
\end{definition*}

\begin{note}
  And indeed, we can still reference them correctly. For instance, we can use \hyperref[def:rational]{rational numbers} to define \hyperref[def:real]{real numbers} and then further use it to define \hyperref[def:complex]{complex numbers}.
\end{note}

Furthermore, we can completely control the name of our environments. We already saw we can name definition, lemma, proposition, corollary and theorem environment. In fact, we can also name remark, note, example and proof as follows.
\begin{eg}[Interesting Example]\label{eg}
  We note that \(1 \neq 2\)!
  \begin{note}[Important note]
    As a consequence, \(2 \neq 3\) also.
  \end{note}

  \begin{remark}[Easy observation]
    We see that from here, we easily have the following theorem.
    \begin{theorem}[Lebesgue Differentiation Theorem]\label{thm:Lebesgue-differentiation-theorem}
      Let \(f\in L^1\), then
      \[
        \lim\limits_{r \to 0} \frac{1}{m(B(x, r))}\int_{B(x, r)}\left\vert f(y) - f(x) \right\vert   \,\mathrm{d}y = 0
      \]
      for a.e. \(x\).
    \end{theorem}
    \begin{proof}[An obvious proof of \autoref{thm:Lebesgue-differentiation-theorem}]
      Obvious.
    \end{proof}
  \end{remark}
\end{eg}
As we can see, specifically for the \texttt{proof} environment, we allow \texttt{autoref} and \texttt{hyperref}. One can actually allow all example, note and remark environment's name to use reference, but I think that is overkilled. But this can be achieved by modify the header in an obvious way.\footnote{This time I mean it!}



\chapter{Known Bugs}
\lecture{2}{9 Sep. 08:00}{Second Lecture}

\section{Introduction}
Nothing is bugs-free. There are some known bugs which I don't have incentive to solve, or it is hard to solve whatsoever. Let me list some of them.

\subsection{Footnote Environment}
It's easy to let you fall into a situation that you want to keep using \texttt{footnote} to add a bunch of unrelated stuffs. However, with our environment there is a known strange behavior, which is following.
\begin{eg}
	Footnote!\footnote{This is a footnote!}
	\begin{remark}
		Oops! footnote somehow shows up earlier than expect!\footnote{This is another footnote!}
	\end{remark}

	Bugs caught!\footnote{The final footnote which is ok!}
\end{eg}

As we saw, the footnote in the \texttt{Example} environment should show at the bottom of its own box, but it's caught by \texttt{Remark} which causes the unwanted behavior. Unfortunately, I haven't found a nice way to solve this. A potential way to solve this is by using \texttt{footnotemark} with \texttt{footnotetext} placing at the bottom of the environment, but this is tedious and needs lots of manual tweaking.

Furthermore, not sure whether you notice it or not, but the color box of \texttt{Remark} is not quite right! It extends to the right, another trick bug...

\subsection{Mdframe Environment}
Though \texttt{mdframe} package is nice and is the key theme throughout this template, but it has some kind of weird behavior. Let's see the demo.

\begin{proof}[Proof of \autoref{thm}]
	We need to prove the followings.
	\begin{claim}
		\(E = mc^{2} \).
	\end{claim}
	\begin{explanation}

		Nonsense.

		Nonsense,

		Nonsense,

		Nonsense,

		Nonsense,

		Nonsense.
	\end{explanation}

\end{proof}

I expect it should break much earlier, and this seems to be an \href{https://github.com/marcodaniel/mdframed/issues/9}{algorithmic issue} of \href{https://ctan.org/pkg/mdframed?lang=en}{\texttt{mdframe}}. One potential solution is to use \href{https://ctan.org/pkg/tcolorbox?lang=en}{\texttt{tcolorbox}} instead, but I haven't completely figure it out, hence I can't really say anything right now.

%─────Lectures{start{end}─────────────────────────────────────────────────────────────────────────────────────────────────────────────────────────────
\lec{1}{9999}

\newpage
%─────Appendix────────────────────────────────────────────────────────────────────────────────────────────────────────────────────────────────────────
%\appendix
%\appendixpage

%\input{appendix.tex}

\newpage
%─────Reference───────────────────────────────────────────────────────────────────────────────────────────────────────────────────────────────────────
%\printbibliography

\end{document}